% GENERATED FILE. Edit canonical lessons, then run npm run generate:books.
% canonical-source-hash: fnv1a64:9da526d3e29d5858
% canonical-lessons: MR-C517-sanketsabda, MR-C517-antarjal, MR-C517-sandes, MR-C517-duva, MR-C517-sathavne

\chapter{Online Life}
\label{ch:mr-online-life}
\hlchaptermodality{517}

\begin{chapteropening}
\textbf{By the end of this chapter:} \emph{I can talk about passwords, messages and links.}

Complete the last lesson of chapter 517: I can talk about passwords, messages and links.
\end{chapteropening}

\section[\texorpdfstring{\mr{संकेतशब्द} (saṅketśabda)}{संकेतशब्द (saṅketśabda)}]{\mr{संकेतशब्द} (saṅketśabda) — a password}

\label{lesson:MR-C517-sanketsabda}



\begin{quote}
Before the new one: say the Marathi for a video clip, then the Marathi for to use.
\end{quote}

\subsection*{You'll want to know: \mr{संकेतशब्द}}
\textbf{\mr{संकेतशब्द}} — \emph{saṅketśabda} — \textquotedblleft{}a password\textquotedblright{}.

\begin{quote}
\mr{मी} \mr{माझा} \mr{संकेतशब्द} \mr{कोणालाही} \mr{सांगत} \mr{नाही}.
\end{quote}

\begin{quote}
\emph{mī mājhā saṅketśabda koṇālāhī sāṅgat nāhī.}
\end{quote}

\begin{quote}
I do not tell anyone my password.
\end{quote}

Masculine.

\begin{cousinweb}
\textbf{\mr{संकेतशब्द}} (\emph{saṅketśabda}) is a tatsama compound of Sanskrit \emph{saṁketa}, \textquotedblleft{}an agreed sign\textquotedblright{}, and \emph{śabda}, \textquotedblleft{}word\textquotedblright{} (Monier-Williams).
\end{cousinweb}

\begin{grammarlens}[title={What you've built}]
One more word for describing what you lived and felt.
\end{grammarlens}

\subsection*{Your turn}
\begin{itemize}
\raggedright
  \item \emph{Say it:} \emph{saṅketśabda}
  \item \emph{Say it:} \emph{saṅketśabda}, once more
  \item \emph{Say it:} \emph{mī mājhā saṅketśabda koṇālāhī sāṅgat nāhī.} — I do not tell anyone my password.
  \item \emph{Recall:} say the Marathi for a video clip, then the Marathi for to use, then say \emph{saṅketśabda} again
\end{itemize}

\begin{tcolorbox}[breakable,colback=teal!4,colframe=teal!35!black,title={Before you move on}]
What does \mr{संकेतशब्द} mean? (\textquotedblleft{}A password\textquotedblright{}.) Say it once more.
\end{tcolorbox}
\section[\texorpdfstring{\mr{आंतरजाल} (āntarjāl)}{आंतरजाल (āntarjāl)}]{\mr{आंतरजाल} (āntarjāl) — the internet}

\label{lesson:MR-C517-antarjal}



\begin{quote}
Before the new one: say the Marathi for to use, then the Marathi for a password.
\end{quote}

\subsection*{You'll want to know: \mr{आंतरजाल}}
\textbf{\mr{आंतरजाल}} — \emph{āntarjāl} — \textquotedblleft{}the internet\textquotedblright{}.

\begin{quote}
\mr{गावात} \mr{आता} \mr{आंतरजाल} \mr{चांगलं} \mr{चालतं}.
\end{quote}

\begin{quote}
\emph{gāvāt ātā āntarjāl cāṅglaṁ cāltaṁ.}
\end{quote}

\begin{quote}
The internet works well in the village now.
\end{quote}

Neuter.

In speech, \mr{इंटरनेट}; \mr{आंतरजाल} is the Marathi word, a net between.

\begin{cousinweb}
\textbf{\mr{आंतरजाल}} (\emph{āntarjāl}) is a modern tatsama coinage from Sanskrit \emph{antara}, \textquotedblleft{}between\textquotedblright{}, and \emph{jāla}, \textquotedblleft{}net\textquotedblright{} (Monier-Williams).
\end{cousinweb}

\begin{grammarlens}[title={What you've built}]
One more word for describing what you lived and felt.
\end{grammarlens}

\subsection*{Your turn}
\begin{itemize}
\raggedright
  \item \emph{Say it:} \emph{āntarjāl}
  \item \emph{Say it:} \emph{āntarjāl}, once more
  \item \emph{Say it:} \emph{gāvāt ātā āntarjāl cāṅglaṁ cāltaṁ.} — The internet works well in the village now.
  \item \emph{Recall:} say the Marathi for to use, then the Marathi for a password, then say \emph{āntarjāl} again
\end{itemize}

\begin{tcolorbox}[breakable,colback=teal!4,colframe=teal!35!black,title={Before you move on}]
What does \mr{आंतरजाल} mean? (\textquotedblleft{}The internet\textquotedblright{}.) Say it once more.
\end{tcolorbox}
\section[\texorpdfstring{\mr{संदेश} (sandeś)}{संदेश (sandeś)}]{\mr{संदेश} (sandeś) — a message}

\label{lesson:MR-C517-sandes}



\begin{quote}
Before the new one: say the Marathi for a password, then the Marathi for the internet.
\end{quote}

\subsection*{You'll want to know: \mr{संदेश}}
\textbf{\mr{संदेश}} — \emph{sandeś} — \textquotedblleft{}a message\textquotedblright{}.

\begin{quote}
\mr{अरुणने} \mr{मित्राला} \mr{एक} \mr{छोटा} \mr{संदेश} \mr{पाठवला}.
\end{quote}

\begin{quote}
\emph{aruṇne mitrālā ek choṭā sandeś pāṭhavlā.}
\end{quote}

\begin{quote}
Arun sent his friend a short message.
\end{quote}

Masculine.

\mr{संदेश} \mr{पाठवणे}, to send a message.

\begin{cousinweb}
\textbf{\mr{संदेश}} (\emph{sandeś}) is a tatsama, taken from Sanskrit \emph{saṁdeśa}, \textquotedblleft{}message, information\textquotedblright{} (Monier-Williams).
\end{cousinweb}

\begin{grammarlens}[title={What you've built}]
One more word for describing what you lived and felt.
\end{grammarlens}

\subsection*{Your turn}
\begin{itemize}
\raggedright
  \item \emph{Say it:} \emph{sandeś}
  \item \emph{Say it:} \emph{sandeś}, once more
  \item \emph{Say it:} \emph{aruṇne mitrālā ek choṭā sandeś pāṭhavlā.} — Arun sent his friend a short message.
  \item \emph{Recall:} say the Marathi for a password, then the Marathi for the internet, then say \emph{sandeś} again
\end{itemize}

\begin{tcolorbox}[breakable,colback=teal!4,colframe=teal!35!black,title={Before you move on}]
What does \mr{संदेश} mean? (\textquotedblleft{}A message\textquotedblright{}.) Say it once more.
\end{tcolorbox}
\section[\texorpdfstring{\mr{दुवा} (duvā)}{दुवा (duvā)}]{\mr{दुवा} (duvā) — a link; a connecting link}

\label{lesson:MR-C517-duva}



\begin{quote}
Before the new one: say the Marathi for the internet, then the Marathi for a message.
\end{quote}

\subsection*{You'll want to know: \mr{दुवा}}
\textbf{\mr{दुवा}} — \emph{duvā} — \textquotedblleft{}a link; a connecting link\textquotedblright{}.

\begin{quote}
\mr{शिक्षकांनी} \mr{सहलीच्या} \mr{चित्रफितीचा} \mr{दुवा} \mr{पाठवला}.
\end{quote}

\begin{quote}
\emph{śikṣakānnī sahlīcyā citraphitīcā duvā pāṭhavlā.}
\end{quote}

\begin{quote}
The teacher sent the link to the video of the trip.
\end{quote}

Masculine.

Once a link of a chain; now also a link on a web page.

\begin{cousinweb}
\textbf{\mr{दुवा}} (\emph{duvā}) is the Marathi word for a link, of a chain or of a web page; its origin is not settled.
\end{cousinweb}

\begin{grammarlens}[title={What you've built}]
One more word for describing what you lived and felt.
\end{grammarlens}

\subsection*{Your turn}
\begin{itemize}
\raggedright
  \item \emph{Say it:} \emph{duvā}
  \item \emph{Say it:} \emph{duvā}, once more
  \item \emph{Say it:} \emph{śikṣakānnī sahlīcyā citraphitīcā duvā pāṭhavlā.} — The teacher sent the link to the video of the trip.
  \item \emph{Recall:} say the Marathi for the internet, then the Marathi for a message, then say \emph{duvā} again
\end{itemize}

\begin{tcolorbox}[breakable,colback=teal!4,colframe=teal!35!black,title={Before you move on}]
What does \mr{दुवा} mean? (\textquotedblleft{}A link; a connecting link\textquotedblright{}.) Say it once more.
\end{tcolorbox}
\section[\texorpdfstring{\mr{साठवणे} (sāṭhavṇe)}{साठवणे (sāṭhavṇe)}]{\mr{साठवणे} (sāṭhavṇe) — to store; to save up}

\label{lesson:MR-C517-sathavne}



\begin{quote}
Before the new one: say the Marathi for a message, then the Marathi for a link.
\end{quote}

\subsection*{You'll want to know: \mr{साठवणे}}
\textbf{\mr{साठवणे}} — \emph{sāṭhavṇe} — \textquotedblleft{}to store; to save up\textquotedblright{}.

\begin{quote}
\mr{मी} \mr{सगळी} \mr{छायाचित्रं} \mr{संगणकात} \mr{साठवली}.
\end{quote}

\begin{quote}
\emph{mī sagḷī chāyācitraṁ saṅgaṇakāt sāṭhavlī.}
\end{quote}

\begin{quote}
I stored all the photos on the computer.
\end{quote}

Past, with -\mr{ने} (he …): \mr{त्याने} \mr{पाणी} \mr{साठवलं}, \emph{tyāne pāṇī sāṭhavlaṁ}.

Where things go takes -\mr{त}: \mr{संगणकात} \mr{साठवणे}.

\begin{cousinweb}
\textbf{\mr{साठवणे}} (\emph{sāṭhavṇe}) is a tadbhava: CDIAL lists Marathi \emph{sā̃ṭhaviṇẽ}, \textquotedblleft{}to heap up\textquotedblright{}, under Sanskrit \emph{saṁsthāpayati}, \textquotedblleft{}sets up, establishes\textquotedblright{} (CDIAL 13051).
\end{cousinweb}

\begin{grammarlens}[title={What you've built}]
One more word for describing what you lived and felt.
\end{grammarlens}

\subsection*{Your turn}
\begin{itemize}
\raggedright
  \item \emph{Say it:} \emph{sāṭhavṇe}
  \item \emph{Say it:} \emph{sāṭhavṇe}, once more
  \item \emph{Say it:} \emph{mī sagḷī chāyācitraṁ saṅgaṇakāt sāṭhavlī.} — I stored all the photos on the computer.
  \item \emph{Recall:} say the Marathi for a message, then the Marathi for a link, then say \emph{sāṭhavṇe} again
\end{itemize}

\begin{tcolorbox}[breakable,colback=teal!4,colframe=teal!35!black,title={Before you move on}]
What does \mr{साठवणे} mean? (\textquotedblleft{}To store; to save up\textquotedblright{}.) Say it once more.
\end{tcolorbox}
